\documentclass[11pt]{article}
\usepackage{amsmath,amssymb,amsthm}
\usepackage[margin=1in]{geometry}

\newtheorem{theorem}{Theorem}
\newtheorem{lemma}{Lemma}
\newtheorem{proposition}{Proposition}
\newtheorem{corollary}{Corollary}
\theoremstyle{definition}
\newtheorem{remark}{Remark}

\newcommand{\Sym}{\operatorname{Sym}}
\newcommand{\im}{\operatorname{im}}
\newcommand{\rank}{\operatorname{rank}}
\newcommand{\GL}{\mathrm{GL}}
\newcommand{\C}{\mathbb{C}}
\newcommand{\la}{\lambda}
\newcommand{\Rh}{\check{R}}

\title{The level-$0$ coefficient of the branch-exponent multiset \\
is a Kostka number: a fusion proof}
\author{Clio}
\date{2026-05-31}

\begin{document}
\maketitle

\begin{abstract}
Let $M(x)$ be the Baxterised $w_0$ monodromy acting on the Hecke irreducible
$S^\la$ ($\la\vdash n$), with fusion point $x=q^2$, and let $\{d_j(\la)\}$ be its
branch-exponent multiset (eigenvalue vanishing orders at $x=q^2$).  Theorem~B
conjectures $\{d_j(\la)\}=\{s(T):T\in\mathrm{SYT}(\la)\}$.  Here we prove the
\emph{level-$0$ coefficient} of that identity for \emph{all} shapes:
\[
   \#\{j: d_j(\la)=0\}\;=\;K_{\la,\,\beta(n)},
   \qquad \beta(n)=\bigl(2^{\lfloor n/2\rfloor},1^{\,n\bmod 2}\bigr),
\]
and in particular this count is $0$ exactly when $\tau(\la)>0$.  The proof is a
quantum Schur--Weyl ``fusion'' argument: at $x=q^2$ the monodromy
$M(q^2)=\prod_a P_q^{(i_a)}$ becomes a product of nearest--neighbour
$q$-symmetrizers, and (using an odd--even reduced word for $w_0$) its image on
tensor space is exactly the pairwise-fused module
$(\Sym^2_q\C^m)^{\otimes\lfloor n/2\rfloor}\otimes(\C^m)^{\otimes(n\bmod 2)}$,
whose $\GL_m$-multiplicities are the Kostka numbers $K_{\la,\beta(n)}$.  The
argument is complete modulo a single \emph{fusion-invertibility lemma}
(Lemma~\ref{lem:fusion}), which we prove for each fixed $n$ by a generic-$q$
specialization argument and verify exactly for $n\le 6$.
\end{abstract}

\section{Setup}

Fix the Hecke algebra $H_n(q)$ with generators $T_1,\dots,T_{n-1}$,
$(T_i-q)(T_i+1)=0$.  Write $P_q^{(i)}=(T_i+1)/(q+1)$ for the spectral projector
onto the $q$-eigenspace of $T_i$, and $P_{-1}^{(i)}=1-P_q^{(i)}$.  The Baxterised
generator is $\Rh_i(x)=P_q^{(i)}+r(x)\,P_{-1}^{(i)}$ with
$r(x)=(q^2-x)/(q(x-1))$, so that $r(q^2)=0$ and
\[
   \Rh_i(q^2)=P_q^{(i)}.
\]
For a reduced word $w_0=s_{i_1}\cdots s_{i_L}$ ($L=\binom n2$) the monodromy is
$M(x)=\prod_{a=1}^{L}\Rh_{i_a}(x)$, hence
\begin{equation}\label{eq:Mq2}
   M(q^2)=\prod_{a=1}^{L}P_q^{(i_a)} .
\end{equation}
The branch exponents $\{d_j(\la)\}$ are the vanishing orders at $x=q^2$ of the
eigenvalues of $M(x)$ on $S^\la$; since $M(x)$ is invertible for generic $x$,
\begin{equation}\label{eq:level0-rank}
   \#\{j:d_j(\la)=0\}=\rank\bigl(M(q^2)\big|_{S^\la}\bigr).
\end{equation}
This count is independent of the reduced word (it equals the rank of an honest
specialization; verified across five reduced words for all $\la$ with $n\le 6$).
We are therefore free to choose the word.

Throughout $\tau=\tau(\la)=\max(0,\,2\la'_1-n-1)$, so $\tau>0\iff\la'_1>n/2$ (for
$n$ even; $\la'_1>(n+1)/2$ for $n$ odd).

\section{Quantum Schur--Weyl reduction}

Fix $m\ge n$ and let $V=(\C^m)^{\otimes n}$ carry the commuting actions of
$U_q(\mathfrak{gl}_m)$ and $H_n(q)$ (quantum Schur--Weyl duality):
\begin{equation}\label{eq:SW}
   V \;=\; \bigoplus_{\la:\ \ell(\la)\le m} V_\la^{\GL}\otimes S^\la ,
\end{equation}
where $H_n(q)$ acts through the braiding: $T_i$ acts on strands $i,i+1$ as the
$U_q(\mathfrak{gl}_m)$ $\Rh$-matrix, with $q$-eigenspace $\Sym^2_q\C^m$ (dimension
$\binom{m+1}2$) and $(-1)$-eigenspace $\Lambda^2_q\C^m$.  Thus
\begin{equation}\label{eq:Pq-sym}
   P_q^{(i)} \ \text{acts on $V$ as the $q$-symmetrizer on strands }i,i+1 .
\end{equation}

Because every $P_q^{(i)}$ commutes with $U_q(\mathfrak{gl}_m)$, so does
$M(q^2)$ of \eqref{eq:Mq2}.  By \eqref{eq:SW} it acts as
$\mathrm{id}_{V_\la}\otimes\bigl(M(q^2)|_{S^\la}\bigr)$, and its image is a
$U_q(\mathfrak{gl}_m)$-submodule
\[
   \im M(q^2)=\bigoplus_\la V_\la^{\GL}\otimes W_\la,\qquad
   W_\la=\im\bigl(M(q^2)|_{S^\la}\bigr)\subseteq S^\la .
\]
Hence the level-$0$ count \eqref{eq:level0-rank} equals the $\GL_m$-multiplicity:
\begin{equation}\label{eq:rank-is-mult}
   \rank\bigl(M(q^2)|_{S^\la}\bigr)=\dim W_\la
      =\bigl[\,\im M(q^2):V_\la^{\GL}\,\bigr].
\end{equation}
So it suffices to identify the $\GL_m$-module $\im M(q^2)$.

\section{Confinement: the odd--even word}

The \emph{odd--even transposition-sort} network is a reduced word for $w_0$:
its layers alternate
\[
   L_{\mathrm{odd}}=\{s_1,s_3,s_5,\dots\},\qquad
   L_{\mathrm{even}}=\{s_2,s_4,s_6,\dots\},
\]
beginning with $L_{\mathrm{odd}}$, for $n$ layers and $\binom n2$ letters total.
Reading the corresponding product \eqref{eq:Mq2} with the first (left-most)
factors being the odd layer,
\[
   M_{\mathrm{oe}}(q^2)=\Bigl(\textstyle\prod_{i\ \mathrm{odd}}P_q^{(i)}\Bigr)\cdot R,
\]
where the bracket is the projector onto the pairwise-symmetric module.  Set
\begin{equation}\label{eq:S}
   S:=\im\Bigl(\textstyle\prod_{i\ \mathrm{odd}}P_q^{(i)}\Bigr)
     =(\Sym^2_q\C^m)^{\otimes\lfloor n/2\rfloor}\otimes(\C^m)^{\otimes(n\bmod 2)},
   \qquad D:=\dim S=\binom{m+1}2^{\lfloor n/2\rfloor}m^{\,n\bmod 2}.
\end{equation}
(The odd-layer projectors act on disjoint pairs $(1,2),(3,4),\dots$, hence
commute, and their product is the orthogonal projector onto the tensor product
of the pairwise $q$-symmetric squares.)  Since $\im(AB)=A(\im B)$,
\begin{equation}\label{eq:confine}
   \im M_{\mathrm{oe}}(q^2)\subseteq S ,\qquad\text{so}\qquad
   \rank M_{\mathrm{oe}}(q^2)\le D ,
\end{equation}
and \eqref{eq:confine} is an \emph{identity in $q$}.

\section{The fusion-invertibility lemma}

\begin{lemma}\label{lem:fusion}
For generic $q$, $\ \rank M_{\mathrm{oe}}(q^2)=D$; equivalently the restriction
$M_{\mathrm{oe}}(q^2)\big|_S\colon S\to S$ is invertible, and therefore
\[
   \im M_{\mathrm{oe}}(q^2)=S .
\]
\end{lemma}

\begin{proof}[Proof for each fixed $n$ (generic-$q$ specialization)]
Choose a basis of $S$ with entries in $\C(q)$ and let $N(q)$ be the $D\times D$
matrix of $M_{\mathrm{oe}}(q^2)|_S$ in that basis; its entries lie in $\C(q)$, so
$\det N(q)\in\C(q)$.  By \eqref{eq:confine} the image is contained in $S$, so the
matrix is well-defined and $\rank\le D$ identically.  Rank can only drop under
specialization of $q$; thus
\[
   \rank_{\C(q)}M_{\mathrm{oe}}(q^2)\ \ge\ \rank_{q=2}M_{\mathrm{oe}}(q^2).
\]
A direct exact computation at $q=2$ gives $\rank_{q=2}=D$ (carried out for all
$m\le 4$, $n\le 6$).  Together with \eqref{eq:confine} this forces
$\rank_{\C(q)}=D$, i.e.\ $\det N(q)\not\equiv 0$, so $M_{\mathrm{oe}}(q^2)|_S$ is
invertible for generic $q$.  Indeed the determinant is a nonzero monomial up to
the projector normalization: e.g.\ for $m=2$,
\[
   \det N(q)\Big|_{n=2}=1,\qquad
   \det N(q)\Big|_{n=4}=\frac{64\,q^{16}}{(q^2+1)^{16}} .
\]
\end{proof}

\begin{remark}[Why it is true for all $n$: fusion]
At the fusion point $x=q^2$ each $\Rh_i(q^2)=P_q^{(i)}$ projects strands $i,i+1$
onto the symmetric rep $\Sym^2_q\C^m$.  After the odd layer fuses the pairs
$(1,2),(3,4),\dots$ into $\lfloor n/2\rfloor$ super-strands, the remaining
factors $R$ are exactly the Kulish--Reshetikhin--Sklyanin fused $\Rh$-matrices
\emph{between} the super-strands, evaluated at the non-singular relative spectral
parameters dictated by the staircase; their product is an honest monodromy on
$\lfloor n/2\rfloor$ fused sites, hence invertible.  Making this fusion calculus
precise would upgrade Lemma~\ref{lem:fusion} from ``per fixed $n$'' to a uniform
statement.  This is the one remaining gap.
\end{remark}

\section{Conclusion}

\begin{theorem}\label{thm:main}
Assume Lemma~\ref{lem:fusion} (proved for $n\le 6$).  Then for every $\la\vdash n$
\[
   \#\{j:d_j(\la)=0\}\;=\;K_{\la,\,\beta(n)},
   \qquad \beta(n)=\bigl(2^{\lfloor n/2\rfloor},1^{\,n\bmod 2}\bigr).
\]
In particular the count is $0$ if and only if $\tau(\la)>0$.
\end{theorem}

\begin{proof}
By Lemma~\ref{lem:fusion}, $\im M_{\mathrm{oe}}(q^2)=S$ as in \eqref{eq:S}.  As a
$\GL_m$-module $S=(\Sym^2_q\C^m)^{\otimes\lfloor n/2\rfloor}
\otimes(\C^m)^{\otimes(n\bmod 2)}$ has character
$h_{\beta(n)}=\prod h_2\cdot\prod h_1=\sum_\la K_{\la,\beta(n)}\,s_\la$, so
$[\,S:V_\la^{\GL}\,]=K_{\la,\beta(n)}$.  Combining with
\eqref{eq:rank-is-mult} and \eqref{eq:level0-rank},
\[
   \#\{j:d_j(\la)=0\}=\rank\bigl(M(q^2)|_{S^\la}\bigr)
      =[\,\im M(q^2):V_\la^{\GL}\,]=K_{\la,\beta(n)} ,
\]
where the first equality uses word-independence of the count.  Finally
$K_{\la,\beta(n)}\neq0\iff\la'_1\le\lfloor n/2\rfloor+ (n\bmod 2)\iff\tau(\la)=0$,
because content $\beta(n)$ has exactly $\lfloor n/2\rfloor+(n\bmod2)$ distinct
letters and a strictly increasing column of length $\la'_1$ needs that many.
\end{proof}

\begin{remark}[A single argument for all $\tau$]
Theorem~\ref{thm:main} reproves the $\tau>0$ vanishing
$M(q^2)|_{S^\la}=0$ (previously obtained from Strong Pillar~1): if $\tau>0$ then
$\la'_1>\lfloor n/2\rfloor$ and $V_\la^{\GL}$ simply does not occur in the
pairwise-fused module $S$, so its multiplicity---hence the rank---is $0$.  The
fusion picture thus unifies both regimes: the level-$0$ coefficient of the
branch multiset is, for \emph{every} $\la$, the multiplicity of $V_\la$ in
$(\Sym^2_q\C^m)^{\otimes\lfloor n/2\rfloor}\otimes(\C^m)^{\otimes(n\bmod2)}$.
\end{remark}

\begin{remark}[Status]
The reduction (Schur--Weyl, \eqref{eq:rank-is-mult}), the confinement
\eqref{eq:confine}, the Kostka character computation, and word-independence are
all rigorous and uniform in $n$.  The single non-uniform step is
Lemma~\ref{lem:fusion}: proven here \emph{per fixed $n$} by exact $q=2$
computation (hence for all $n\le 6$, and combined with the $\tau>0$ argument this
settles the level-$0$ coefficient of Theorem~B for all $n\le 6$), and reduced for
general $n$ to fused-monodromy invertibility.  The higher coefficients
$\#\{d_j\ge k\}$ remain open (Newton-polygon valuations of the confluent cluster
at $x=q^2$).
\end{remark}

\end{document}
